\section{Metodología experimental}

\subsection{Características de la máquina}

Los experimentos se realizarán en una estación de trabajo con las siguientes
características:
\begin{itemize}
    \item \textbf{Procesador:} Intel Core i5-13420H (Raptor Lake), 8 núcleos y
    12 hilos (4 P-Cores y 4 E-Cores), frecuencia base de 2.10 GHz y turbo de
    hasta 4.60 GHz. Cachés: L1d de 320 KiB, L1i de 384 KiB, L2 de 7 MiB y L3
    compartida de 12 MiB.
    \item \textbf{Memoria principal:} 16 GiB DDR (15 GiB utilizables por el
    sistema operativo).
    \item \textbf{Almacenamiento:} SSD NVMe WD PC SN740 de 512 GB, interfaz
    PCIe Gen4 $\times$4.
    \item \textbf{Entorno de compilación:} Arch Linux x86\_64 (kernel 6.x),
    GCC y las banderas \texttt{-std=c++20 -O3 -march=native}.
    \item \textbf{Afinidad de CPU:} se fijará el proceso de medición al núcleo
    lógico 2 mediante \texttt{taskset -c 2}. Esto reduce las migraciones entre
    núcleos, pero no elimina variaciones de frecuencia, temperatura ni la
    actividad concurrente del sistema.
\end{itemize}

\subsection{Datos de entrada}

Se utilizarán \textbf{arreglos de enteros generados artificialmente}.
Para cada tamaño \(n\) se generarán tres
arreglos con valores de distribución uniforme en el intervalo
\([0,10^9]\), utilizando las semillas fijas \textbf{1729, 2718 y 31415}.
Se registrarán las semillas y los parámetros de generación para reproducir
las entradas. Sparse Table y Segment Tree recibirán exactamente los mismos
datos en cada comparación.

\begin{itemize}
    \item \textbf{Representación de los elementos:} cada entero se asociará
    con su posición mediante un par \((\text{valor},\text{índice})\).
    Ambos componentes serán enteros: el dato numérico y su posición.
    El orden lexicográfico permitirá obtener el índice del mínimo y,
    ante valores iguales, seleccionar la posición más a la izquierda.
    La preparación de estos pares quedará fuera del tiempo medido.
    \item \textbf{Consultas:} para una longitud de rango \(k\), con
    \(1\leq k\leq n\), se elegirá uniformemente una posición inicial
    \(L\) entre \(0\) y \(n-k\), y se fijará \(R=L+k-1\).
    Así, cada consulta corresponderá a un intervalo inclusivo válido
    \([L,R]\), que se traducirá a \([L,R+1)\) al invocar las
    implementaciones. Se considerarán distintas longitudes y ambas
    estructuras responderán la misma secuencia de consultas.
    \item \textbf{Actualizaciones:} en el escenario dinámico se generará
    una secuencia reproducible de posiciones elegidas uniformemente entre
    \(0\) y \(n-1\), permitiendo posiciones repetidas, y nuevos valores
    del intervalo \([0,10^9]\), distintos del valor vigente.
    Los cambios se aplicarán en el mismo orden: Segment Tree actualizará
    la hoja correspondiente y Sparse Table se reconstruirá a partir del
    arreglo modificado. El índice asociado a cada elemento se conservará.
    Las consultas posteriores serán iguales para ambas alternativas.
    \item \textbf{Entradas de validación:} se emplearán, por separado,
    arreglos pequeños crecientes, decrecientes, con todos sus elementos
    iguales, con mínimos repetidos y de un solo elemento. Estos casos
    comprobarán la corrección y el desempate; no constituirán experimentos
    temporales adicionales.
\end{itemize}

La generación de los arreglos, las consultas y la secuencia de modificaciones
se realizará antes del cronometraje. Los tamaños de arreglo, las longitudes
concretas de los rangos y las cantidades de operaciones se fijarán antes de
ejecutar los experimentos.

\subsection{Experimentos e hipótesis}

Se compararán Sparse Table y Segment Tree mediante \textbf{cuatro experimentos}:
construcción inicial y consultas en el escenario estático, y mantenimiento
y consultas posteriores a modificaciones en el escenario dinámico. Para cada
entrada se construirán ambas estructuras y se \textbf{reutilizarán las
instancias} para consultar el arreglo sin cambios. Después se aplicará la
secuencia de modificaciones: Segment Tree conservará su instancia mediante
\texttt{update}, mientras que Sparse Table será sustituida por una nueva
construcción del arreglo vigente. No se reconstruirán estructuras entre
consultas si no ha ocurrido un cambio.

El cronometraje seguirá el esquema de la plantilla \texttt{uhr}, de Leonardo
Lovera (LELE), adaptado para usar \texttt{std::chrono::steady\_clock} y
expresar los tiempos en nanosegundos \cite{lovera-uhr}. Se realizarán
\textbf{32 repeticiones por configuración y semilla}, después del
calentamiento. Las consultas estáticas se medirán en lotes calibrados fuera
de muestra, conservando el tamaño del lote entre repeticiones de una misma
configuración y método. En el escenario dinámico se respetará la carga fija
definida a continuación, sin añadir consultas para calibrar el reloj.
El tiempo de cada lote se dividirá por la cantidad de consultas ejecutadas.
La construcción y el mantenimiento se expresarán por construcción y por
modificación, respectivamente, sin mezclarlos con el tiempo de consulta.

La generación de datos, preparación de pares, validación y escritura de
resultados quedarán fuera del reloj. También se excluirá la destrucción de
las estructuras, tanto después de construir como al reemplazar Sparse Table;
en este último caso se cronometrará la creación de la tabla nueva. Los índices
devueltos se consumirán para evitar que el compilador elimine las consultas.
La corrección se contrastará con un barrido lineal independiente, fuera de
los intervalos medidos.

En cada repetición dinámica se partirá del mismo arreglo inicial y se
reproducirán los mismos cambios y consultas. La restauración del estado y el
calentamiento quedarán fuera de la muestra; dentro de la secuencia se
conservarán las instancias y su estado entre operaciones. Los experimentos
de mantenimiento y consultas dinámicas se distinguirán por sus intervalos
de cronometraje, sin duplicar una modificación en el mismo recorrido.

Se ejecutarán \textbf{10 ciclos por arreglo, semilla y repetición dinámica}.
Cada ciclo comenzará con \textbf{una actualización puntual}, seguida de
\textbf{1.000 consultas}: primero se aplicará \texttt{update} en Segment
Tree o se reconstruirá Sparse Table y, solo después, se consultará el estado
vigente. Esto supone \textbf{10 modificaciones y 10.000 consultas medidas por
estructura en cada repetición dinámica}. Las 1.000 consultas de cada ciclo
se repartirán equilibradamente entre las longitudes de rango seleccionadas,
registrando por separado el tiempo y la cantidad de consultas de cada
longitud. Así se reutilizará cada estado actualizado para todos sus lotes,
sin multiplicar las reconstrucciones al cambiar la longitud del rango.

No se forzará que las consultas aleatorias incluyan la posición modificada,
para conservar la distribución de rangos definida. Como comprobación de
corrección, después del lote se verificará también la posición actualizada
mediante una consulta unitaria fuera del cronometraje.

Se conservarán muestras crudas y resúmenes CSV con \textbf{media, desviación
estándar muestral, mínimo, Q1, mediana, Q3 y máximo}, identificando estructura,
tamaño del arreglo, longitud del rango, semilla, repetición y ciclo cuando
corresponda. Para resumir el escenario dinámico se obtendrá, en cada
repetición, el tiempo medio por modificación sobre sus diez ciclos y el
tiempo por consulta de cada longitud sobre sus lotes. Los ciclos no se
contarán como repeticiones independientes adicionales. Las figuras
mostrarán media y barras de una desviación estándar, como en el boletín
anterior; no se interpretarán como intervalos de confianza. Cada semilla
conservará su resumen: las repeticiones sobre un mismo arreglo no se tratarán
como arreglos independientes. Las diferencias se evaluarán junto con su
dispersión: las 32 repeticiones permiten caracterizar la variabilidad, pero
no garantizan por sí solas significancia estadística ni una superioridad
universal a partir de las cotas asintóticas.

Antes de la ejecución oficial se realizará un \textbf{piloto}, especialmente
con el arreglo de mayor tamaño, para estimar la duración de la matriz
completa. El presupuesto del conjunto será de \textbf{menos de una hora},
con una proyección objetivo de \textbf{45 minutos como máximo} y
\textbf{15 minutos de margen}. La estimación incluirá calentamiento,
calibración, restauración y destrucción de estructuras, validación, análisis
de memoria, escritura y generación de gráficos, aunque estas tareas estén
fuera del cronometraje de las operaciones. Si la proyección supera el
objetivo, se revisará el tamaño máximo o la cantidad de ciclos dinámicos
antes de ejecutar, manteniendo las 32 repeticiones y documentando el ajuste.
No se dará por garantizada la duración hasta comprobarla con el piloto.

\subsubsection{Construcción inicial en el escenario estático}

\begin{itemize}
    \item \textbf{Variable:} \(n\), cantidad de elementos del arreglo.
    \item \textbf{Controles:} mismos valores y pares de entrada para ambas
    estructuras, con las tres semillas definidas. Cada construcción se
    realizará desde la entrada preparada, sin incluir su generación.
    \item \textbf{Resultados:} tiempo por construcción frente a \(n\),
    mediante tablas y gráficos. Las instancias construidas se reutilizarán
    para las consultas; su costo inicial no se cargará a cada consulta.
\end{itemize}

\textbf{Hipótesis.} Se espera que Sparse Table tenga mayor tiempo de
construcción al crecer \(n\), porque prepara una tabla de tamaño
\(\Theta(n\log n)\), mientras que Segment Tree construye sus nodos en
\(\Theta(n)\). El preprocesamiento adicional de Sparse Table puede
compensarse cuando se realizan suficientes consultas estáticas, pero este
beneficio se estudiará separando ambos costos, no suponiendo que ocurre para
cualquier cantidad de consultas.

\subsubsection{Consultas en el escenario estático}

\begin{itemize}
    \item \textbf{Variables:} tamaño del arreglo \(n\) y longitud del
    rango \(k\). Se variará una manteniendo fija la otra en cada comparación.
    \item \textbf{Controles:} instancias ya construidas, sin modificaciones
    del arreglo, y mismos lotes de intervalos aleatorios. Se comparará
    \texttt{query\_idempotent} de Sparse Table con \texttt{query} de
    Segment Tree.
    \item \textbf{Resultados:} tiempo por consulta frente a \(n\) para
    longitudes fijas y frente a \(k\) para tamaños de arreglo fijos.
    La construcción inicial quedará excluida de estos tiempos.
\end{itemize}

\textbf{Hipótesis.} Se espera que Sparse Table tenga menor tiempo por
consulta al resolver RMQ con dos accesos y una combinación, de costo
\(\Theta(1)\), aprovechando su preprocesamiento. Segment Tree recorrerá
los niveles necesarios para descomponer el rango, con peor caso
\(\Theta(\log n)\). No se presupone una diferencia apreciable para todos
los rangos: los unitarios son constantes en ambas estructuras, y con \(k\)
fijo la cota del recorrido del árbol no obliga a crecer con \(n\).

\subsubsection{Mantenimiento en el escenario dinámico}

\begin{itemize}
    \item \textbf{Variable:} tamaño del arreglo \(n\), bajo modificaciones
    puntuales de sus valores.
    \item \textbf{Controles:} mismo estado inicial y misma secuencia de
    posiciones y nuevos valores. Se realizarán diez ciclos, cada uno con
    una modificación seguida de 1.000 consultas compartidas entre las
    longitudes evaluadas. Cada cambio deberá quedar reflejado antes de
    responderlas; no se agruparán modificaciones para reconstruir
    Sparse Table una sola vez.
    \item \textbf{Resultados:} tiempo por modificación frente a \(n\).
    Se medirá \texttt{update} en Segment Tree y la reconstrucción en
    Sparse Table, excluyendo las consultas y la construcción inicial.
\end{itemize}

\textbf{Hipótesis.} Se espera que Segment Tree tenga menor tiempo de
mantenimiento, ya que actualiza una hoja y sus ancestros en
\(\Theta(\log n)\), frente a reconstruir Sparse Table en
\(\Theta(n\log n)\). La diferencia debería aumentar al crecer el
arreglo. Esta comparación expresa el costo de mantener respuestas correctas
tras un cambio, no dos operaciones nativas de actualización equivalentes.

\subsubsection{Consultas en el escenario dinámico}

\begin{itemize}
    \item \textbf{Variables:} tamaño del arreglo \(n\) y longitud del
    rango \(k\), manteniendo fija una al estudiar la otra.
    \item \textbf{Controles:} mismos cambios previos y lotes de consultas.
    Antes de consultar, Segment Tree deberá estar actualizado y Sparse Table
    reconstruida, de modo que ambas representen el mismo arreglo vigente.
    Se usarán los lotes de 1.000 consultas de los diez ciclos del experimento
    de mantenimiento, sin repetir los cambios ni las reconstrucciones para
    obtener estas mediciones.
    \item \textbf{Resultados:} tiempo por consulta después de las
    modificaciones, separado del tiempo de mantenimiento registrado en el
    experimento anterior. Las respuestas deberán coincidir con el mínimo
    y el desempate definidos.
\end{itemize}

\textbf{Hipótesis.} Una vez mantenidas las estructuras, Sparse Table debería
conservar su ventaja en consultas por su costo \(\Theta(1)\); Segment Tree
mantiene su cota de peor caso \(\Theta(\log n)\). Por tanto, la ventaja
esperada de Segment Tree en un entorno con cambios se atribuye al menor costo
de mantenimiento, no necesariamente a una consulta aislada más rápida.
Los tiempos pueden diferir de los estáticos por el estado de la caché dejado
por las actualizaciones o reconstrucciones, sin que cambien estas cotas.

\subsubsection{Análisis del espacio utilizado}

Se analizarán las asignaciones de los constructores para los tamaños de
arreglo evaluados, considerando el mismo tipo de par en ambas estructuras.
Se distinguirá el almacenamiento retenido de las reservas temporales de
construcción o reconstrucción, excluyendo el arreglo externo compartido.
Cualquier seguimiento de asignaciones se realizará en una pasada separada,
sin activarlo durante el cronometraje; estos bytes no representan la memoria
residente total del proceso.

\textbf{Hipótesis.} Sparse Table requerirá más espacio al crecer \(n\),
por sus \(n(\lfloor\log_2 n\rfloor+1)\) elementos almacenados, frente
a los \(2n\) de Segment Tree. Las actualizaciones del árbol reutilizan el
vector existente, mientras que la reconstrucción de la tabla requiere nuevas
reservas y puede aumentar el pico si coexisten las versiones anterior y nueva.
Este análisis acompañará los cuatro experimentos temporales, sin constituir
un quinto experimento de tiempo.

\textit{Parámetros por concretar antes de la ejecución: tamaños de arreglos,
longitudes de rangos, cantidad de consultas estáticas y criterios de
calibración de sus lotes. Los diez ciclos dinámicos propuestos quedarán
sujetos a la validación del presupuesto mediante el piloto, sin reducir
las 32 repeticiones.}
